\section*{الفصل \ref{ch:b2:structure-determination} --- تحديد البنية: الرنين المغناطيسي النووي للكربون \texorpdfstring{\protect\babelsublr{$^{13}$C}}{13C} \omterm{def:b2:mass-spec-atomic:spectrum}{ومطيافية الكتلة} والأشعة تحت الحمراء معًا}

\begin{solution}{exo:b2:structure-determination:1}
1،2-ثنائي ميثيل بنزين: 4 (\ce{CH3}؛ C1/C2؛ C3/C6؛ C4/C5). و1،3-: 5 (\ce{CH3}؛ C1/C3؛ C2؛ C4/C6؛ C5).
و1،4-: 3 (\ce{CH3}؛ C1/C4؛ C2/C3/C5/C6).
\end{solution}

\begin{solution}{exo:b2:structure-determination:2}
\ce{CH3CH(OH)CH2CH3}. في DEPT-135: C1 (\ce{CH3}) إلى الأعلى، وC2 (CH) إلى الأعلى، وC3 (\ce{CH2}) إلى الأسفل، وC4 (\ce{CH3}) إلى الأعلى.
وفي DEPT-90: C2 وحده.
\end{solution}

\begin{solution}{exo:b2:structure-determination:3}
209: كربونيل كيتون؛ 170: كربونيل إستر (أو حمض)؛ 129: كربون \omterm{def:b2:huckel:aromatic}{عطري}؛ 64: كربون مرتبط
بأكسجين؛ 14: \ce{CH3} ألكيلية.
\end{solution}

\begin{solution}{exo:b2:structure-determination:4}
$(2 \times 8 + 2 - 8)/2 = 5$: حلقة بنزين (4) و\ce{C=O} (1). مثلًا بنزوات الميثيل،
\ce{C6H5COOCH3}، أو حمض 4-ميثيل بنزويك.
\end{solution}

\begin{solution}{exo:b2:structure-determination:5}
درجة عدم تشبع واحدة، وكيتون (\qty{1715}{cm^{-1}}، والخط عند 209، رباعي)؛ ومجموعة \ce{CH2} واحدة ومجموعتا
\ce{CH3}، وأربع ذرات كربون مختلفة: البيوتان-2-ون، \ce{CH3COCH2CH3}.
\end{solution}

\begin{solution}{exo:b2:structure-determination:6}
الكربون: يُظهر البنتان-2-ون خمسة خطوط، والبنتان-3-ون المتناظر ثلاثة. والبروتونات: للبنتان-2-ون إشارة
أحادية لثلاثة بروتونات (\ce{CH3CO})؛ وللبنتان-3-ون رباعية وثلاثية فقط. \omterm{def:b2:mass-spec-atomic:spectrum}{وطيف الكتلة}: للبنتان-2-ون
قمة أساس عند 43 (\ce{CH3CO+}) وأيون ماكلافرتي عند 58؛ وللبنتان-3-ون قمة أساس عند 57
(\ce{C2H5CO+}).
\end{solution}

\begin{solution}{exo:b2:structure-determination:7}
\ce{C4H8O2}، 88: زوج الرباعية--الثلاثية هو \ce{OCH2CH3} (فمجموعة \ce{CH2} عند 4.12 على الأكسجين)، 
والأحادية هي \ce{CH3CO}. إنه إيثانوات الإيثيل، \ce{CH3COOCH2CH3}.
\end{solution}

\begin{solution}{exo:b2:structure-determination:8}
1،2-ثنائي كلوروبنزين: ثلاثة خطوط (C1/C2، C3/C6، C4/C5). و1،4-ثنائي كلوروبنزين: خطان (C1/C4، 
وCH الأربع). ويُظهر DEPT-90 خطي CH للأول، وخطًا واحدًا للثاني.
\end{solution}

\begin{solution}{exo:b2:structure-determination:9}
C2 مركز فراغي. واستبدال الديوتيريوم بأحد بروتوني C3 أو بالآخر يعطي
متماكبين لامرآويين، لا مرآويين: فالبروتونان لامرآويا الموضع، في محيطين مختلفين.
فيمكن أن يكون لهما انزياحان مختلفان وأن يقترن أحدهما بالآخر وبجيرانهما معًا، فتعطي
\ce{CH2} متعددة معقدة بدلًا من خماسية بسيطة.
\end{solution}

\begin{solution}{exo:b2:structure-determination:10}
خمس درجات عدم تشبع؛ وحلقة أحادية الاستبدال (ثلاثة خطوط CH وخط C واحد عند 137.0، وخمسة
بروتونات \omterm{def:b2:huckel:aromatic}{عطرية})؛ وكيتون (200.8، C)؛ ومجموعة إيثيل (\ce{CH2} 31.8 رباعية عند 2.98 مجاورة للكربونيل، و\ce{CH3}
8.3). إنه 1-فينيل بروبان-1-ون، \ce{C6H5COCH2CH3}. والترافق مع الحلقة يجعل إلكترونات $\pi$ في \ce{C=O}
غير متمركزة ويضعف الرابطة: فتنخفض الحزمة تحت \qty{1715}{cm^{-1}} الخاصة بالكيتون المشبع.
% ledger: ir:CO-ketone
\end{solution}

\begin{solution}{exo:b2:structure-determination:11}
حمض البنتانويك: حزمة \ce{O-H} العريضة جدًا من 2500 إلى \qty{3300}{cm^{-1}} وبروتون قرب 11--12
ppm. وبيوتانوات الميثيل: إشارة أحادية لثلاثة بروتونات في مجال البروتونات المحمولة على كربون مرتبط بأكسجين
(3.3--4.5). وبروبانوات الإيثيل: رباعيتان، إحداهما على الأكسجين والأخرى مجاورة للكربونيل. وإيثانوات
البروبيل: إشارة أحادية لثلاثة بروتونات مجاورة للكربونيل (2.0--2.4) وثلاثية على الأكسجين.
% ledger: ir:OH-acid, nmrr:acid, nmrr:alcohol, nmrr:ketone
\end{solution}

\begin{solution}{exo:b2:structure-determination:12}
يختفي الخطان كلاهما في \omterm{def:b2:structure-determination:dept}{DEPT}، فكلاهما رباعي، ولا يستطيع \omterm{def:b2:structure-determination:dept}{DEPT} وحده التمييز بينهما. أما الانزياحات
فتستطيع: فكربونيلات الإسترات تقع قرب 167--171 والكربونات \omterm{def:b2:huckel:aromatic}{العطرية} بين 128 و137 في مخطط هذا
الفصل. وكربون حلقي عند 166.5 وكربونيل عند 130.6 سيقعان كلاهما بعيدًا خارج صنفيهما؛ 
والنسب الصحيح هو العكس.
\end{solution}

\begin{solution}{pb:b2:structure-determination:1}
\textbf{1.} $3.1/31.7 \approx 9.8~\%$؛ $9.8/1.11 \approx 9$ ذرات كربون.
\textbf{2.} كتلة زوجية: لا نيتروجين (أو ذرتان).
\textbf{3.} $150 - 108 = 42$: يناسب \ce{H10O2} (10 + 32): \ce{C9H10O2}.
\textbf{4.} عشر ذرات كربون كانت ستعطي $M+1 \approx 11~\%$ من $M$، لا 9.8~\%؛ والأشعة تحت الحمراء تُظهر
كربونيلًا، يتطلب ذرة أكسجين ثانية إلى جانب حزمة \ce{C-O}.
\textbf{5.} $(2 \times 9 + 2 - 10)/2 = 5$.
\textbf{6.} $108 + 10 \times 1.007825 + 2 \times 15.994915 \approx 150.0681$.
\textbf{7.} اهتزاز تمدد \ce{C=O} عند موضع إستر مشبع (قرب \qty{1735}{cm^{-1}}).
\textbf{8.} اهتزاز تمدد \ce{C-O} للإستر.
\textbf{9.} روابط \ce{C-H} \omterm{def:b2:huckel:aromatic}{عطرية} (أو ألكينية).
\textbf{10.} مجموعة \ce{O-H}: فلا كحول، ولا حمض كربوكسيلي.
\textbf{11.} لا: فالكربونيل المترافق مع حلقة يمتص عند عدد موجي أدنى؛ وهو هنا حيث تقع
الإسترات المشبعة، فثمة مجموعة غير مترافقة تفصله عن الحلقة.
\textbf{12.} 170.70: \ce{C=O} الإستر؛ 136.14: كربون الحلقة الحامل للبديل؛ 128.56 و128.24:
CH \omterm{def:b2:huckel:aromatic}{عطرية}؛ 66.24: \ce{CH2} مرتبطة بأكسجين؛ 20.82: \ce{CH3}.
\textbf{13.} 66.24، أي \ce{CH2}.
\textbf{14.} 170.70 و136.14، أي ذرتا الكربون اللتان لا هيدروجين عليهما.
\textbf{15.} 128.56 و128.24.
\textbf{16.} أربعة: الكربون الحامل للبديل، وذرتا أورثو، وذرتا ميتا وCH في الموضع بارا. فالمتوقع ثلاثة
أنواع من CH؛ واثنان منها انزياحاهما أقرب من أن يُميَّزا عند هذا المجال، فيتقاسمان خطًا واحدًا.
\textbf{17.} ليس بالضرورة: فالارتفاعات ليست أعدادًا (\cref{prop:b2:structure-determination:intensities})،
وإن كان الخط هنا قد يضم فعلًا نوعين متراكبين من CH.
\textbf{18.} 7.33: البروتونات \omterm{def:b2:huckel:aromatic}{العطرية} الخمسة؛ 5.09: \ce{OCH2}؛ 2.06: \ce{CH3CO}.
\textbf{19.} لأن الذرات المجاورة لها (كربون الحلقة والأكسجين بالنسبة إلى \ce{CH2}، وكربون الكربونيل
بالنسبة إلى \ce{CH3}) لا تحمل هيدروجينًا.
\textbf{20.} إنها مرتبطة بأكسجين الإستر الكهرسلبي وبالحلقة: وكلاهما يقلل حجبها.
\textbf{21.} \ce{C6H5CH2OCOCH3}، إيثانوات البنزيل.
\textbf{22.} في المتماكب لا تكون \ce{CH2} مرتبطة بأكسجين، فلا يمكن أن تظهر عند 5.09، 
والميثيل، الواقع على الأكسجين، كان سيظهر في المجال 3.3--4.5، لا عند 2.06 بجوار كربونيل.
% ledger: nmrr:alcohol, nmrr:ketone
\textbf{23.} 108: فقد 42 (الكيتين، \ce{CH2=C=O}) بإعادة ترتيب، أي الكاتيون الجذري للصيغة
\ce{C7H8O}؛ 91: \ce{C7H7+}، كاتيون البنزيل؛ 43: \ce{CH3CO+}.
\textbf{24.} سبعة أنواع من الكربون (\ce{C=O}، وكربون الحلقة المستبدَل، وCH في المواضع أورثو وميتا وبارا،
و\ce{CH2}، و\ce{CH3}): \textbf{سبعة خطوط، يتجه أحدها (\ce{CH2}) إلى الأسفل} في DEPT-135.
\end{solution}
